\scanpage{196}
\mspage{5.}
Notons d'ailleurs

\uline{Prop. 4} Pour tt schéma \deleted{à opérateurs} $Z$ sur \uline{$S$}, à groupe fini $G$ d'opérateurs, \deleted{\unsure{opé}} agissant de façon admissible, le champ \deleted{défini}\inserted{décrit} \deleted{par la pr}\inserted{comme} cidessus est un champ en \inserted{$\mathcal U$-}catégories multigaloisiennes, et le foncteur ``passage au quotient par $G$'' (qui est défini, car l'action de $G$ reste admissible) \deleted{est} \unsure{font|fait} une donnée de ramification sur $S$.

Tout $\mathscr R_{S'}$ est multigaloisien : tout d'abord, il en est ainsi de la catégorie des \deleted{objets} rev. à opérateurs $G$ \deleted{\unsure{du|des}} $Z_{S'}$, or on en prend ici une sous-catégorie stable par $\mathop{\mathrm{lim}}\limits_{\substack{\longrightarrow\\\longleftarrow}}$ finies et par sommandes directs ; O.K. Foncteurs de chgt de base exact : O.K.

Le foncteur $r$, étant un passage au quotient, commute aux $\varinjlim$ finies (qui ici se calculent au sens de $(Sch)$ !), d'où R1. De $=$ R2 est immédiat. R3 aussi. R4 aussi. O.K.

\deleted{Ainsi} Supposons pour simplifier $S$ \unsure{qu cpct|qu. cpct}, et que $\mathscr R$ satisfasse une petite condition supplémentaire (tjrs vérifiée en pratique) : \deleted{\unsure{tout obj}}

Pour tt objet de $\mathscr R_S$, son rang sur l'objet final de $\mathscr R_S$ \deleted{est}\inserted{reste} borné, i.e $X$ est subordonné à un \deleted{objet} objet à opérateurs galoisien $(P,G)$ de $\mathscr R_S$.

Alors $\mathscr R_S$ est limite inductive \uline{filtrante} de ses sous-catégories de la forme $\mathscr R_S^P$. On a donc une assez bonne description de $\mathscr R_S$, grâce au théorème, en termes d'un système \uline{projectif} d'objets à opérateurs $(Z_i;G_i)$.
% archive sheet number: 195
